% Candidate figures for main-text inclusion.
% Each is labeled C1--C8 matching the visual-opportunity analysis.

\subsection{Candidate Figures for Main-Text Inclusion}
\label{sec:appendix:candidates}

The eight figures below correspond to visual opportunities identified during manuscript
review.  Each is annotated with the section it would supplement and the unique insight
it contributes beyond the existing prose.  They are presented together here so all
candidates can be evaluated before individual placement decisions are made.

% ============================================================
% C1: Clock Domain Divergence
% ============================================================

\subsubsection*{C1 --- Clock Domain Divergence and Invocation-Order Matching}
\noindent\textit{Supplements \secref{sec:clockdomain}.  Shows why timestamp-based
joining fails and how ordinal-position matching solves it.}

\begin{figure}[H]
\centering
\begin{tikzpicture}[font=\small, >=Stealth]

  %--- nsys timeline ---
  \node[anchor=east, font=\small\bfseries] at (-0.1,3.25) {\nsys};
  \node[anchor=east, font=\scriptsize, gray] at (-0.1,2.95) {CUPTI GPU clock};
  \draw[->,thick] (0,3.1) -- (11.8,3.1);

  % nsys kernels: K_A, K_B, K_C, K_D, K_E (all unique)
  \draw[fill=blue!20,draw=blue!50] (0.55,2.85) rectangle (1.45,3.35);
  \node[font=\scriptsize] at (1.0,3.1) {$K_A$};

  \draw[fill=green!20,draw=green!60!black] (2.05,2.85) rectangle (2.95,3.35);
  \node[font=\scriptsize] at (2.5,3.1) {$K_B$};

  \draw[fill=cyan!20,draw=cyan!60!black] (3.55,2.85) rectangle (4.45,3.35);
  \node[font=\scriptsize] at (4.0,3.1) {$K_C$};

  \draw[fill=orange!20,draw=orange!60] (5.55,2.85) rectangle (6.45,3.35);
  \node[font=\scriptsize] at (6.0,3.1) {$K_D$};

  \draw[fill=red!15,draw=red!50] (7.55,2.85) rectangle (8.45,3.35);
  \node[font=\scriptsize] at (8.0,3.1) {$K_E$};

  %--- ncu timeline ---
  \node[anchor=east, font=\small\bfseries] at (-0.1,0.9) {\ncu};
  \node[anchor=east, font=\scriptsize, gray] at (-0.1,0.6) {replay clock};
  \draw[->,thick] (0,0.75) -- (11.8,0.75);

  % ncu kernels: same names, shifted right by accumulating drift
  % K_A: +0.3, K_B: +0.6, K_C: +0.9, K_D: +1.5, K_E: +2.0
  \draw[fill=blue!20,draw=blue!50] (0.85,0.5) rectangle (1.75,1.0);
  \node[font=\scriptsize] at (1.3,0.75) {$K_A$};

  \draw[fill=green!20,draw=green!60!black] (2.65,0.5) rectangle (3.55,1.0);
  \node[font=\scriptsize] at (3.1,0.75) {$K_B$};

  \draw[fill=cyan!20,draw=cyan!60!black] (4.45,0.5) rectangle (5.35,1.0);
  \node[font=\scriptsize] at (4.9,0.75) {$K_C$};

  \draw[fill=orange!20,draw=orange!60] (7.05,0.5) rectangle (7.95,1.0);
  \node[font=\scriptsize] at (7.5,0.75) {$K_D$};

  \draw[fill=red!15,draw=red!50] (9.55,0.5) rectangle (10.45,1.0);
  \node[font=\scriptsize] at (10.0,0.75) {$K_E$};

  %--- correct ordinal joins (green) ---
  \draw[green!60!black,thick] (1.0,2.85) -- (1.3,1.0);
  \draw[green!60!black,thick] (2.5,2.85) -- (3.1,1.0);
  \draw[green!60!black,thick] (4.0,2.85) -- (4.9,1.0);
  \draw[green!60!black,thick] (6.0,2.85) -- (7.5,1.0);
  \draw[green!60!black,thick] (8.0,2.85) -- (10.0,1.0);

  %--- wrong timestamp join (red dashed) ---
  % nsys K_D at x=6.0; ncu K_C at x=4.9 is closer than ncu K_D at x=7.5
  \draw[dashed,red!75,very thick] (6.0,2.85) -- (4.9,1.0);
  \node[font=\scriptsize,red!75,align=center] at (4.9,1.95)
    {\textit{wrong:}\\$K_D \!\leftrightarrow\! K_C$};

  %--- divergence bracket ---
  \draw[<->,orange!80!black,thick] (6.0,2.0) -- (7.5,2.0);
  \node[above,font=\scriptsize,orange!80!black] at (6.75,2.0)
    {$\Delta t$ grows per kernel};

  %--- legend ---
  \draw[green!60!black,thick] (9.0,2.1) -- (9.6,2.1);
  \node[right,font=\scriptsize] at (9.6,2.1) {IOM: ordinal join};
  \draw[dashed,red!75,thick] (9.0,1.75) -- (9.6,1.75);
  \node[right,font=\scriptsize] at (9.6,1.75) {timestamp join (fails)};

\end{tikzpicture}
\caption{\textbf{C1: Clock domain divergence and invocation-order matching.}
  The \nsys CUPTI GPU clock and the \ncu injection replay clock diverge by
  10--100\,$\mu$s per kernel due to PMU serialization overhead; the gap
  \emph{accumulates} across the trace.  A naïve timestamp join (dashed red)
  maps the \nsys record of $K_D$ to the \ncu row for $K_C$---a wrong counter
  assignment.  Invocation-Order Matching (green) connects each kernel in the
  \nsys trace to the matching kernel in the \ncu output by name and position,
  making the join clock-independent.}
\label{fig:cand:clockdomain}
\end{figure}

% ============================================================
% C2: Attribution Priority Decision Tree
% ============================================================

\subsubsection*{C2 --- Attribution Priority Decision Tree}
\noindent\textit{Supplements \secref{sec:tprofiler}--\secref{sec:heuristic}.
  Shows the hierarchical decision logic for the three attribution paths; not
  expressible as a table.}

\begin{figure}[H]
\centering
\begin{tikzpicture}[
  font=\small, >=Stealth,
  decision/.style={diamond, draw, fill=yellow!15, align=center,
                   inner sep=2pt, aspect=2, font=\small},
  result/.style={rounded corners=4pt, draw, align=center,
                 minimum width=4.2cm, inner sep=5pt},
  high/.style={result, fill=blue!15, draw=blue!50!black},
  med/.style={result, fill=cyan!12, draw=cyan!50!black},
  unattr/.style={result, fill=gray!20, draw=gray!50},
  arr/.style={-Stealth, thick},
]

  % Input kernel
  \node[draw, rounded corners=3pt, fill=white, font=\small\bfseries,
        inner sep=5pt] (k) at (0,0) {kernel $k$\\(name, invocation index)};

  % Decision 1: corr.json
  \node[decision] (d1) at (0,-2.2)
    {\texttt{corr.json} entry\\for \texttt{(name, idx)}?};
  \draw[arr] (k) -- (d1);

  % HIGH result
  \node[high] (high) at (7,-2.2)
    {\textbf{HIGH}\\torch.profiler CUPTI\\correlation};
  \draw[arr] (d1) -- node[above,font=\scriptsize]{YES} (high);

  % Decision 2: NVTX enclosure
  \node[decision] (d2) at (0,-5.0)
    {NVTX enclosure at\\host launch time $T_k$?};
  \draw[arr] (d1) -- node[right,font=\scriptsize]{NO} (d2);

  % MEDIUM NVTX
  \node[med] (mnvtx) at (7,-5.0)
    {\textbf{MEDIUM}\\NVTX enclosure\\(innermost range)};
  \draw[arr] (d2) -- node[above,font=\scriptsize]{YES} (mnvtx);

  % Decision 3: Inductor map
  \node[decision] (d3) at (0,-7.8)
    {Inductor fusion map\\entry for kernel name?};
  \draw[arr] (d2) -- node[right,font=\scriptsize]{NO} (d3);

  % MEDIUM Inductor
  \node[med] (mind) at (7,-7.8)
    {\textbf{MEDIUM}\\Inductor fusion\\enrichment};
  \draw[arr] (d3) -- node[above,font=\scriptsize]{YES} (mind);

  % UNATTRIBUTED
  \node[unattr] (unattr) at (0,-10.2)
    {\textbf{UNATTRIBUTED}\\cuDNN-internal; eager Triton\\without NVTX or debug output};
  \draw[arr] (d3) -- node[right,font=\scriptsize]{NO} (unattr);

  % Augmentation side-note (applies to all tiers; arrow from MEDIUM Inductor for clarity)
  \node[draw=orange!60, fill=orange!8, rounded corners=3pt,
        align=left, font=\scriptsize, inner sep=5pt] (aug) at (7,-10.2)
    {\textit{Augmentation pass (all tiers):}\\
     Inductor fusion map also checked\\
     after attribution is assigned;\\
     \texttt{fused\_ops} list is populated\\
     without changing the confidence tier.};
  \draw[dashed, orange!60] (mind.south) -- (aug.north);

\end{tikzpicture}
\caption{\textbf{C2: Attribution priority decision tree.}
  Each kernel is tested against the three attribution paths in descending confidence
  order.  A HIGH-confidence entry from the \texttt{torch.profiler} correlation pass
  (\secref{sec:tprofiler}) pre-empts the NVTX and Inductor checks.  The Inductor
  fusion enrichment pass (\secref{sec:heuristic}) is applied as a \emph{post-attribution}
  augmentation at all tiers: it populates \texttt{fused\_ops} without overriding
  an existing attribution.  Only kernels that fail all three paths land in
  \texttt{unattributed\_kernels[]}.}
\label{fig:cand:decisiontree}
\end{figure}

% ============================================================
% C4: Three-Level FX Pass Architecture
% ============================================================

\subsubsection*{C4 --- Three-Level FX Pass Hook Points in the \texttt{torch.compile()} Pipeline}
\noindent\textit{Supplements \secref{sec:threelevel}.  Shows \emph{why} the Level~1/2/3
  separation is necessary by locating each hook in the IR lowering chain.}

\begin{figure}[H]
\centering
\begin{tikzpicture}[font=\small, >=Stealth,
  stage/.style={draw, rounded corners=4pt, fill=white,
                minimum width=3.6cm, minimum height=0.75cm,
                align=center, inner sep=5pt, font=\small},
  hook/.style={draw, rounded corners=3pt, align=left,
               inner sep=5pt, font=\scriptsize},
  arr/.style={-Stealth,thick},
  harr/.style={-Stealth, thick, dashed},
]

  % Pipeline stages (vertical, left column)
  \node[stage, fill=gray!10] (in)     at (0,0)    {\texttt{workload.py}};
  \node[stage, fill=blue!8]  (dynamo) at (0,-2.2)  {Dynamo FX\\graph capture};
  \node[stage, fill=blue!8]  (aot)    at (0,-4.4)  {AOTAutograd\\decomposition};
  \node[stage, fill=blue!8]  (ind)    at (0,-6.6)  {Inductor\\code generation};
  \node[stage, fill=gray!10] (out)    at (0,-8.8)  {Triton / CUDA kernels};

  \draw[arr] (in)     -- (dynamo);
  \draw[arr] (dynamo) -- (aot);
  \draw[arr] (aot)    -- (ind);
  \draw[arr] (ind)    -- (out);

  % Level 1 hook (between Dynamo and AOT)
  \node[hook, fill=red!8, draw=red!50] (l1) at (5.5,-2.2)
    {\textbf{Level 1} (pre-decomposition)\\
     \textbullet\ QKV fusion (\textit{SDPA case study})\\
     \textbullet\ Attention canonicalization};
  \draw[harr, red!60] (l1.west) -- (dynamo.east);

  % Level 2 hook (between AOT and Inductor)
  \node[hook, fill=orange!8, draw=orange!50] (l2) at (5.5,-4.4)
    {\textbf{Level 2} (Aten IR)\\
     \textbullet\ BF16 dtype promotion (\textit{all case studies})\\
     \textbullet\ Weight folding, layout annotation};
  \draw[harr, orange!60] (l2.west) -- (aot.east);

  % Level 3 hook (at Inductor)
  \node[hook, fill=green!8, draw=green!50!black] (l3) at (5.5,-6.6)
    {\textbf{Level 3} (Inductor config)\\
     \textbullet\ Weight freezing (\textit{all case studies})\\
     \textbullet\ Autotuning mode selection};
  \draw[harr, green!60!black] (l3.west) -- (ind.east);

  % Compose label
  \node[font=\scriptsize, gray, align=center] at (0,-9.6)
    {Passes at each level compose\\without interfering across levels};

\end{tikzpicture}
\caption{\textbf{C4: Three-level FX pass hook points.}
  The \texttt{torch.compile()} pipeline lowers a model through three IR stages.
  Level~1 passes must run before AOTAutograd because structural patterns such as
  a shared \texttt{LayerNorm} output feeding Q, K, and V projections---the
  precondition for QKV fusion---are decomposed into per-consumer copies by
  AOTAutograd and become invisible afterward.  Level~2 (Aten IR) applies
  dtype and layout transforms on the fully decomposed graph.  Level~3
  Inductor flags control non-graph behaviors such as weight freezing and
  autotuning that cannot be expressed as graph rewrites.}
\label{fig:cand:fxlevel}
\end{figure}


